\section{Discussion}

\subsection{Strategic Style Versus Raw Leaderboard Position}

The benchmark surfaces stable strategic profiles rather than only a leaderboard. Across the frozen player-game aggregate, Kimi has the highest overall win rate at 56.2\% with only 11.3\% challenge frequency, while Mistral remains at 0\% win rate with 63.4\% challenge frequency. Kimi is the most robust finisher because it leads Experiments~1, 2, and~3. Nemotron is the strongest condition-specific finisher in the low-strategy control, where it leads Experiment~0 at 49.0\%. Qwen is the most consistent runner-up across the deception-relevant conditions. These are not interchangeable styles: the benchmark is exposing different ways to trade off bluffing, inspection, and risk.

\begin{figure}[t]
  \centering
  \includegraphics[width=\linewidth]{lie_frequency_vs_win_rate.png}
  \caption{Lie frequency versus win rate across all 24 model-condition rows. Winning is not a monotonic function of deception alone: Kimi combines strong win rates with relatively low challenge intensity, Nemotron is strongest mainly in the control condition, and Mistral occupies the high-suspicion, zero-win corner.}
  \label{fig:lie-frequency-vs-win-rate}
\end{figure}

\begin{figure}[t]
  \centering
  \includegraphics[width=\linewidth]{challenge_frequency_vs_win_rate.png}
  \caption{Challenge frequency versus win rate across all 24 model-condition rows. The strongest outcomes cluster away from the extreme high-challenge region, while Mistral anchors the costly over-enforcement corner.}
  \label{fig:challenge-frequency-vs-win-rate}
\end{figure}

\subsection{Instruction Compliance Under Incentive Pressure}

The honesty-mandate condition is the strongest alignment result in the paper. The plain-language rule reduces deception but does not eliminate it, and it also makes the table less vigilant. This means the main result in Experiment~3 is not only that models still lie. It is that the prompt softens the enforcement regime around those lies. Relative to Experiment~1, mean model-level lie frequency falls from 39.6\% to 27.2\%, but mean challenge frequency also falls from 45.3\% to 21.8\%, while mean lie success rises from 11.0\% to 34.7\%. Prompt framing therefore changes the social equilibrium of the table, not just the actor's willingness to bluff.

\subsection{Inspection, Enforcement, and Calibration}

Bullshit can be read as an inspection game. If players challenge too often, they punish themselves; if they challenge too rarely, bluffing becomes cheap. The frozen cohort supports this interpretation. Experiment~1 produces the highest challenge rate and the lowest lie success, while Experiment~3 produces much lower challenge intensity and much higher lie success. A secondary observation sharpens the point: Experiments~0 and~3 converge to nearly the same enforcement regime, with mean challenge frequencies of 21.5\% and 21.8\% and mean lie-success rates of 37.9\% and 34.7\% respectively. Figures~\ref{fig:lie-frequency-vs-win-rate} and~\ref{fig:challenge-frequency-vs-win-rate} make the same calibration point from two directions: winning is not a simple function of more bluffing or more challenging. In this sense, BS-Bench measures not only who lies, but how agents calibrate signaling, inspection, and punishment under uncertainty.

Mistral is the clearest example of costly over-enforcement. It challenges more than every other model in every condition---62.5\% in Experiment~0, 74.7\% in Experiment~1, 71.3\% in Experiment~2, and 50.8\% in Experiment~3---yet it never wins a condition. This is evidence for a narrow game-theoretic claim about this benchmark: indiscriminate suspicion is a losing strategy. Challenge behavior matters, but only when it is calibrated well enough that the punisher is not paying the cost of constant failed enforcement.

The reasoning traces are consistent with this interpretation. In one Experiment~3 game, Qwen explicitly acknowledges the honesty rule and still chooses a minimal bluff because it has no legal truthful play of the required rank. In another Experiment~3 game, Nemotron challenges an impossible two-card claim by noting that it already holds three of the four relevant cards. Short excerpts are included in Appendix~\ref{app:reasoning-traces} to give a concrete sense of how these strategic choices appear in the logs.

\subsection{Limitations}

The current paper studies one environment family, one primary provider cohort, and one roster of six hosted models. It does not include a human baseline in the main pilot, and it does not support claims about model size effects or broad real-world deception beyond the benchmark. These are limitations of scope, not reasons the benchmark is invalid.
